% ch41_model_categories.tex — Volume IV Chapter 5

\chapter{Model Categories: Axioms and Examples}

\begin{overviewbox}
A \term{model category} is a category equipped with three distinguished
classes of morphisms — \emph{weak equivalences}, \emph{cofibrations}, and
\emph{fibrations} — satisfying axioms that together provide an explicit,
small-hom-class construction of the localization
$\mathcal{C}[\mathcal{W}^{-1}]$ of Chapter~40. The framework is due to
Quillen (1967), and it transformed homotopy theory from an art form into
a systematic apparatus.

The five axioms are: $2$-out-of-$3$ for weak equivalences; retract-closure
for all three classes; \emph{factorization} of every map as cofibration
followed by trivial fibration, or as trivial cofibration followed by
fibration; \emph{lifting} of cofibrations against trivial fibrations
(and dually, trivial cofibrations against fibrations); and completeness
of the underlying category.

This chapter develops the axioms, introduces lifting properties as the
combinatorial heart of the theory, and concludes with the two
foundational examples: the Quillen model structure on simplicial sets
(fibrations $=$ Kan fibrations, cofibrations $=$ monomorphisms) and the
Quillen model structure on topological spaces (fibrations $=$ Serre
fibrations, cofibrations generated by sphere inclusions). The
\term{Quillen equivalence} $\lvert{-}\rvert : \mathbf{sSet} \rightleftarrows
\mathbf{Top} : \mathrm{Sing}$ is stated; its proof — the homotopy
hypothesis at the level of model categories — is the climax of Chapter~42.

Joyal's model structure on $\mathbf{sSet}$, in which fibrant objects are
exactly the quasi-categories, appears at the end of the chapter as a
preview of Chapter~43.
\end{overviewbox}


% ═══════════════════════════════════════════════════════════════════════════
\section{Quillen's Axioms}
\label{sec:quillen-axioms}
% ═══════════════════════════════════════════════════════════════════════════

\subsection{Formalizing}

\begin{definition}[Model category]
\label{def:model-category}
A \term{model category} is a complete and cocomplete category $\mathcal{C}$
together with three classes of morphisms — \term{weak equivalences}
$\mathcal{W}$, \term{cofibrations} $\mathcal{C}\mathrm{of}$, and
\term{fibrations} $\mathcal{F}\mathrm{ib}$ — satisfying the following
five axioms.

\medskip
\textbf{(M1) Two-out-of-three.} If two of $f, g, gf$ are weak
equivalences, so is the third.

\medskip
\textbf{(M2) Retracts.} If $f$ is a retract of $g$ (in the arrow category
of $\mathcal{C}$) and $g$ belongs to $\mathcal{W}$, $\mathcal{C}\mathrm{of}$,
or $\mathcal{F}\mathrm{ib}$, then so does $f$.

\medskip
\textbf{(M3) Factorization.} Every morphism $f$ admits factorizations
\begin{equation*}
  f \;=\; q \circ i \quad \text{with } i \in \mathcal{C}\mathrm{of}, \, q \in \mathcal{W} \cap \mathcal{F}\mathrm{ib}
\end{equation*}
and
\begin{equation*}
  f \;=\; p \circ j \quad \text{with } j \in \mathcal{W} \cap \mathcal{C}\mathrm{of}, \, p \in \mathcal{F}\mathrm{ib}.
\end{equation*}

\medskip
\textbf{(M4) Lifting.} Cofibrations have the left lifting property
against trivial fibrations; trivial cofibrations have the left lifting
property against fibrations. That is, for every commutative square
\begin{equation*}
  \begin{tikzcd}
    A \arrow[r] \arrow[d, "i"'] & X \arrow[d, "p"] \\
    B \arrow[r] \arrow[ur, dashed, "h"] & Y
  \end{tikzcd}
\end{equation*}
with $i \in \mathcal{C}\mathrm{of}$ and $p \in \mathcal{W} \cap
\mathcal{F}\mathrm{ib}$, a lift $h$ exists making both triangles commute.
The same holds if $i \in \mathcal{W} \cap \mathcal{C}\mathrm{of}$ and
$p \in \mathcal{F}\mathrm{ib}$.

\medskip
\textbf{(M5) Completeness.} $\mathcal{C}$ has all small limits and
colimits.
\end{definition}

\begin{howtosay}
Three classes pronounced ``weak equivalences,'' ``cofibrations,'' and
``fibrations.'' \term{Trivial cofibration} $=$ cofibration $\cap$ weak
equivalence (also called \term{acyclic cofibration}); \term{trivial
fibration} $=$ fibration $\cap$ weak equivalence. Common diagrammatic
shorthand: $\xrightarrow{\sim}$ for a weak equivalence, $\hookrightarrow$ for
a cofibration, $\twoheadrightarrow$ for a fibration.
\end{howtosay}

\begin{remark}[Determination of two classes by the third]
\label{rem:two-determine-three}
The five axioms imply that any two of $\mathcal{W}, \mathcal{C}\mathrm{of},
\mathcal{F}\mathrm{ib}$ determine the third:
\begin{itemize}
  \item Cofibrations $=$ maps with LLP against trivial fibrations.
  \item Fibrations $=$ maps with RLP against trivial cofibrations.
  \item Weak equivalences $=$ maps factoring as trivial cofibration $\circ$ trivial fibration (and equivalent characterizations).
\end{itemize}
Hence model structures are often specified by giving only the classes
$\mathcal{W}$ and one other class — the third being determined.
\end{remark}


% ═══════════════════════════════════════════════════════════════════════════
\section{Lifting Properties and Small Objects}
\label{sec:lifting-small-objects}
% ═══════════════════════════════════════════════════════════════════════════

\subsection{Formalizing}

\begin{definition}[Left and right lifting properties]
\label{def:llp-rlp}
A morphism $i : A \to B$ has the \term{left lifting property} (LLP) with
respect to $p : X \to Y$ if every commutative square
\begin{equation*}
  \begin{tikzcd}
    A \arrow[r] \arrow[d, "i"'] & X \arrow[d, "p"] \\
    B \arrow[r] \arrow[ur, dashed] & Y
  \end{tikzcd}
\end{equation*}
admits a diagonal lift. Equivalently, $p$ has the \term{right lifting
property} (RLP) with respect to $i$. For a class $\mathcal{S}$ of
morphisms, $\mathrm{LLP}(\mathcal{S})$ is the class of morphisms with
LLP against every $p \in \mathcal{S}$; dually $\mathrm{RLP}(\mathcal{S})$.
\end{definition}

\begin{lemma}[Closure properties of LLP/RLP]
\label{lem:llp-closure}
$\mathrm{LLP}(\mathcal{S})$ is closed under:
\begin{itemize}
  \item Composition.
  \item Pushouts (in $\mathcal{C}$).
  \item Transfinite composition (sequential colimits along ordinals).
  \item Retracts.
\end{itemize}
$\mathrm{RLP}(\mathcal{S})$ is closed under composition, pullbacks, transfinite limits, and retracts.
\end{lemma}

\begin{proof}[Proof sketch]
For composition: a lift in $A \xrightarrow{i} B \xrightarrow{j} C$ against
$p$ is constructed in two stages, lifting $i$ against $p$ then $j$ against
$p$. For pushouts and transfinite composition: standard diagram-chase,
constructing the lift component-by-component. For retracts: any commutative
square involving the retract is filled in by lifting against the
retracting square. \qed
\end{proof}

\begin{theorem}[Small object argument]
\label{thm:soa}
Let $\mathcal{I}$ be a set of morphisms in a cocomplete category
$\mathcal{C}$ whose domains are \term{small} (every map from such a domain
into a transfinite colimit factors through a finite stage). Then every
morphism $f : X \to Y$ in $\mathcal{C}$ admits a functorial factorization
$f = p \circ i$ with $i \in \mathrm{LLP}(\mathrm{RLP}(\mathcal{I}))$
(i.e., $i$ a cell-complex-like map built from $\mathcal{I}$-cells) and
$p \in \mathrm{RLP}(\mathcal{I})$.
\end{theorem}

\begin{remark}[Why the small object argument matters]
The small object argument is the engine producing factorizations (M3) in
practice. \term{Cofibrantly generated model categories} are those whose
cofibrations and trivial cofibrations are each generated by small sets
$\mathcal{I}$ and $\mathcal{J}$ respectively; the factorizations come
from applying Theorem~\ref{thm:soa} to $\mathcal{I}$ and $\mathcal{J}$.
Both Quillen-$\mathbf{sSet}$ and Quillen-$\mathbf{Top}$ are cofibrantly
generated.
\end{remark}


% ═══════════════════════════════════════════════════════════════════════════
\section{Cofibrant, Fibrant, and Bifibrant Objects}
\label{sec:cofibrant-fibrant}
% ═══════════════════════════════════════════════════════════════════════════

\subsection{Formalizing}

Let $\emptyset$ and $1$ denote the initial and terminal objects of
$\mathcal{C}$ (existing by M5).

\begin{definition}[Cofibrant, fibrant, bifibrant]
\label{def:cofib-fib}
An object $X \in \mathcal{C}$ is:
\begin{itemize}
  \item \term{Cofibrant} if the unique map $\emptyset \to X$ is a cofibration.
  \item \term{Fibrant} if the unique map $X \to 1$ is a fibration.
  \item \term{Bifibrant} if it is both cofibrant and fibrant.
\end{itemize}
\end{definition}

\begin{theorem}[Functorial replacement]
\label{thm:replacement}
There exist endofunctors $Q, R : \mathcal{C} \to \mathcal{C}$ and
natural transformations $Q \xrightarrow{\sim} \mathrm{id}$ and
$\mathrm{id} \xrightarrow{\sim} R$ such that for every $X$:
\begin{itemize}
  \item $Q X$ is cofibrant and $Q X \to X$ is a trivial fibration —
        \term{cofibrant replacement}.
  \item $X \to R X$ is a trivial cofibration and $R X$ is fibrant —
        \term{fibrant replacement}.
\end{itemize}
\end{theorem}

\begin{proof}[Proof sketch]
For cofibrant replacement, factor $\emptyset \to X$ using M3 as
$\emptyset \hookrightarrow Q X \xrightarrow{\sim,\,\twoheadrightarrow} X$.
For fibrant replacement, factor $X \to 1$ as
$X \xrightarrow{\sim,\,\hookrightarrow} R X \twoheadrightarrow 1$.
Functoriality of the factorizations (M3 with functorial choice via the
small object argument) gives the natural transformations. \qed
\end{proof}

\begin{remark}[Why bifibrant?]
The localization $\mathrm{Ho}(\mathcal{C}) = \mathcal{C}[\mathcal{W}^{-1}]$
is built from $\mathcal{C}$'s objects, but its hom-sets between
bifibrant $X, Y$ admit the explicit description
\begin{equation*}
  \mathrm{Ho}(\mathcal{C})(X, Y) \;=\; \mathcal{C}(X, Y) / \sim,
\end{equation*}
where $\sim$ is the equivalence relation of \term{homotopy} (defined below
via cylinder/path objects). Combined with bifibrant replacement (apply
$Q$ then $R$, or any composite producing a bifibrant object), this
gives the small-hom-class construction of $\mathrm{Ho}(\mathcal{C})$.
\end{remark}

\subsection{Reorganizing: cylinders and homotopies}

\begin{definition}[Cylinder object, path object]
\label{def:cylinder-path}
A \term{cylinder object} for $X$ is a factorization $X \sqcup X
\xrightarrow{(i_0, i_1)} \mathrm{Cyl}(X) \xrightarrow{\sim} X$ of the
fold map, with $(i_0, i_1)$ a cofibration. A \term{path object} for $Y$
is a factorization $Y \xrightarrow{\sim} \mathrm{Path}(Y) \xrightarrow{(p_0, p_1)}
Y \times Y$ of the diagonal, with $(p_0, p_1)$ a fibration.
Two maps $f, g : X \to Y$ are \term{left homotopic} if there is $H :
\mathrm{Cyl}(X) \to Y$ with $H i_0 = f$, $H i_1 = g$; \term{right
homotopic} if there is $H : X \to \mathrm{Path}(Y)$ with $p_0 H = f$,
$p_1 H = g$.
\end{definition}

\begin{theorem}[Homotopy is an equivalence relation on bifibrant maps]
\label{thm:homotopy-equiv}
If $X$ is cofibrant and $Y$ is fibrant, left and right homotopy coincide
and define an equivalence relation on $\mathcal{C}(X, Y)$, which we denote
$[X, Y]$. The composition $\mathcal{C}(X, Y) \times \mathcal{C}(Y, Z) \to
\mathcal{C}(X, Z)$ descends to $[-,-]$, giving the homotopy category
$\mathrm{Ho}(\mathcal{C})_{\mathrm{bif}}$ on bifibrant objects.
\end{theorem}


% ═══════════════════════════════════════════════════════════════════════════
\section{The Two Foundational Examples}
\label{sec:foundational-examples}
% ═══════════════════════════════════════════════════════════════════════════

\subsection{$\mathbf{sSet}$ with the Quillen model structure}

\begin{theorem}[Quillen model structure on $\mathbf{sSet}$]
\label{thm:sset-quillen}
The category $\mathbf{sSet}$ admits a model structure in which:
\begin{itemize}
  \item Weak equivalences are maps whose realization is a weak homotopy
        equivalence in $\mathbf{Top}$.
  \item Cofibrations are monomorphisms (equivalently, level-wise
        injections).
  \item Fibrations are \term{Kan fibrations}: maps $p : X \to Y$ with
        RLP against every horn inclusion $\Lambda^n_k \hookrightarrow
        \Delta^n$ ($n \geq 1$, $0 \leq k \leq n$).
\end{itemize}
This is the \term{Kan--Quillen model structure}. Generating cofibrations:
the boundary inclusions $\partial\Delta^n \hookrightarrow \Delta^n$.
Generating trivial cofibrations: the horn inclusions $\Lambda^n_k
\hookrightarrow \Delta^n$.
\end{theorem}

\begin{remark}[Fibrant objects of Quillen-$\mathbf{sSet}$]
The fibrant objects are exactly the Kan complexes (Chapter~39): $X \to 1$
having RLP against all horn inclusions is exactly the Kan condition.
Every simplicial set is cofibrant (the unique map from $\emptyset$ is a
monomorphism), so bifibrant simplicial sets are Kan complexes.
\end{remark}

\subsection{$\mathbf{Top}$ with the Quillen model structure}

\begin{theorem}[Quillen model structure on $\mathbf{Top}$]
\label{thm:top-quillen}
$\mathbf{Top}$ (compactly generated weak Hausdorff) admits a model
structure with:
\begin{itemize}
  \item Weak equivalences: weak homotopy equivalences (isomorphisms on
        all $\pi_n$).
  \item Cofibrations: retracts of relative cell complexes, generated by
        the sphere inclusions $S^{n-1} \hookrightarrow D^n$.
  \item Fibrations: \term{Serre fibrations}: maps with RLP against
        $D^n \hookrightarrow D^n \times [0, 1]$.
\end{itemize}
Every space is fibrant; cofibrant spaces are (retracts of) CW complexes.
\end{theorem}

\subsection{The Quillen equivalence}

\begin{theorem}[Realization–singular Quillen equivalence]
\label{thm:realization-sing-quillen}
The adjunction $\lvert{-}\rvert \dashv \mathrm{Sing}$ of Chapter~39 is a
\term{Quillen equivalence} between Quillen-$\mathbf{sSet}$ and
Quillen-$\mathbf{Top}$: it descends to an equivalence of homotopy
categories
\begin{equation*}
  \mathrm{Ho}(\mathbf{sSet}) \;\simeq\; \mathrm{Ho}(\mathbf{Top}).
\end{equation*}
\end{theorem}

\begin{remark}[The homotopy hypothesis]
Theorem~\ref{thm:realization-sing-quillen} is the formal version of the
\term{homotopy hypothesis}: the homotopy theory of spaces is equivalent
to the homotopy theory of simplicial sets, hence — restricting to
Kan complexes — to the homotopy theory of $\infty$-groupoids.
Chapter~42 develops the proof via derived functors. We will use this
equivalence throughout Volume~IV to move between topological and
simplicial language as convenient.
\end{remark}

\subsection{Preview: the Joyal model structure}

\begin{remark}[Joyal's model structure]
\label{rem:joyal-preview}
$\mathbf{sSet}$ also admits a \emph{different} model structure, due to
Joyal, in which:
\begin{itemize}
  \item Cofibrations are still monomorphisms (same as Kan--Quillen).
  \item Fibrations are \term{quasi-fibrations}: maps with RLP against
        \emph{inner} horn inclusions $\Lambda^n_k \hookrightarrow \Delta^n$
        ($0 < k < n$).
  \item Weak equivalences are characterized by induced equivalences of
        homotopy $1$-categories after fibrant replacement.
\end{itemize}
The fibrant objects are the \term{quasi-categories} (Chapter~43). The
Joyal model structure makes $\mathbf{sSet}$ a model for $\infty$-categories
exactly as Kan--Quillen makes it a model for $\infty$-groupoids.
\end{remark}


% ═══════════════════════════════════════════════════════════════════════════
\section{Exercises}
\label{sec:model-exercises}
% ═══════════════════════════════════════════════════════════════════════════

\begin{exercisesbox}
\begin{qa}
\qaitem{What is a model category?}{A complete and cocomplete category with three classes of morphisms — weak equivalences $\mathcal{W}$, cofibrations $\mathcal{C}\mathrm{of}$, fibrations $\mathcal{F}\mathrm{ib}$ — satisfying $2$-out-of-$3$, retract closure, factorization, lifting, and completeness (M1--M5).}
\qaitem{What is a trivial cofibration?}{A cofibration that is also a weak equivalence. Equivalently: a map with LLP against all fibrations.}
\qaitem{What is a trivial fibration?}{A fibration that is also a weak equivalence. Equivalently: a map with RLP against all cofibrations.}
\qaitem{State the lifting axiom (M4).}{Cofibrations lift on the left against trivial fibrations; trivial cofibrations lift on the left against fibrations.}
\qaitem{State the factorization axiom (M3).}{Every map factors as cofib $\circ$ trivial fib, and also as trivial cofib $\circ$ fib.}
\qaitem{Why is one of $\mathcal{W}, \mathcal{C}\mathrm{of}, \mathcal{F}\mathrm{ib}$ determined by the other two?}{Lifting plus factorization: cofibrations are exactly $\mathrm{LLP}(\mathcal{W} \cap \mathcal{F}\mathrm{ib})$, fibrations are $\mathrm{RLP}(\mathcal{W} \cap \mathcal{C}\mathrm{of})$, and $\mathcal{W}$ is determined by either via the factorizations.}
\qaitem{What is the left lifting property $\mathrm{LLP}(\mathcal{S})$?}{The class of maps $i$ such that every square with $i$ on the left and some $p \in \mathcal{S}$ on the right admits a diagonal lift.}
\qaitem{What closure properties does $\mathrm{LLP}(\mathcal{S})$ enjoy?}{Composition, pushouts, transfinite composition, retracts. These are the building blocks of ``cell complex'' arguments.}
\qaitem{What is the small object argument?}{Given a small set of morphisms $\mathcal{I}$ with small domains, every map factors functorially as (transfinite cell-pushout of $\mathcal{I}$) followed by (map with RLP against $\mathcal{I}$).}
\qaitem{What is a cofibrantly generated model category?}{One where $\mathcal{C}\mathrm{of}$ and $\mathcal{W} \cap \mathcal{C}\mathrm{of}$ are each generated by a small set under closure of $\mathrm{LLP}(\mathrm{RLP}(\cdot))$. Both Quillen-$\mathbf{sSet}$ and Quillen-$\mathbf{Top}$ are cofibrantly generated.}
\qaitem{Define cofibrant and fibrant.}{$X$ cofibrant: $\emptyset \to X$ is a cofibration. $X$ fibrant: $X \to 1$ is a fibration. Bifibrant: both.}
\qaitem{Why do we need replacement?}{To represent arbitrary objects by bifibrant ones, on which the homotopy relation is well-behaved. Factorize $\emptyset \to X$ and $X \to 1$ to get $QX$ and $RX$.}
\qaitem{State the homotopy-category description.}{If $X$ is cofibrant and $Y$ is fibrant, $\mathrm{Ho}(\mathcal{C})(X, Y) = \mathcal{C}(X, Y)/\sim$, where $\sim$ is the (left $=$ right) homotopy equivalence relation.}
\qaitem{What is the Kan--Quillen model structure on $\mathbf{sSet}$?}{Cofibrations $=$ monomorphisms; fibrations $=$ Kan fibrations (RLP against horn inclusions); weak equivalences $=$ realization-weak-homotopy-equivalences. Fibrant objects $=$ Kan complexes.}
\qaitem{What is the Quillen model structure on $\mathbf{Top}$?}{Cofibrations $=$ retracts of relative cell complexes (generated by $S^{n-1} \hookrightarrow D^n$); fibrations $=$ Serre fibrations; weak equivalences $=$ weak homotopy equivalences. Every space is fibrant; CW complexes are cofibrant.}
\qaitem{What is the Quillen equivalence between $\mathbf{sSet}$ and $\mathbf{Top}$?}{The adjunction $\lvert{-}\rvert \dashv \mathrm{Sing}$ descends to an equivalence $\mathrm{Ho}(\mathbf{sSet}) \simeq \mathrm{Ho}(\mathbf{Top})$. This is the homotopy hypothesis.}
\qaitem{What is the Joyal model structure?}{Same cofibrations as Kan--Quillen; fibrations $=$ inner-fibrations (RLP against \emph{inner} horns only); weak equivalences different. Fibrant objects $=$ quasi-categories (Chapter~43).}
\qaitem{Why use cofibrant replacement before computing $\pi_n$?}{Because $\pi_n$ is only well-defined and homotopy-invariant on bifibrant objects. Cofibrant replacement ensures the inputs to representable functors are correctly resolved.}
\qaitem{What's a cylinder object for $X$?}{A factorization $X \sqcup X \hookrightarrow \mathrm{Cyl}(X) \xrightarrow{\sim} X$ of the fold map. Two maps $X \to Y$ are left-homotopic via a $\mathrm{Cyl}(X) \to Y$.}
\qaitem{Path object dually?}{$Y \xrightarrow{\sim} \mathrm{Path}(Y) \twoheadrightarrow Y \times Y$ factoring the diagonal. Right homotopies live here.}
\qaitem{Why does the homotopy category have small hom-sets in model categories?}{Because for bifibrant $X, Y$, hom-sets are quotients of $\mathcal{C}(X, Y)$ (small) by an equivalence relation, hence small. This solves the size issue of Gabriel--Zisman localization (Chapter~40).}
\qaitem{Is $\mathbf{sSet}$ with the Joyal structure a model for ordinary categories?}{No — for $\infty$-categories. Ordinary categories are the homotopy 1-category. Joyal-$\mathbf{sSet}$ retains higher coherence data (Chapter~43).}
\qaitem{What's special about a Kan fibration?}{It satisfies the homotopy lifting property for all simplicial sets, including for outer horns — modeling Serre fibrations of topological spaces.}
\end{qa}
\end{exercisesbox}


% ═══════════════════════════════════════════════════════════════════════════
\section*{Chapter Summary}
\addcontentsline{toc}{section}{Chapter Summary}
% ═══════════════════════════════════════════════════════════════════════════

\begin{summarybox}
\textbf{Model category.}\quad Complete cocomplete category with three classes
$\mathcal{W}, \mathcal{C}\mathrm{of}, \mathcal{F}\mathrm{ib}$ satisfying
$2$-out-of-$3$, retract closure, factorization, lifting, and completeness.
Two classes determine the third via lifting properties.

\textbf{Lifting and the small object argument.}\quad LLP/RLP define the
classes from each other. The small object argument produces factorizations
from generating sets of cofibrations.

\textbf{Cofibrant/fibrant.}\quad $\emptyset \to X$ cof, $X \to 1$ fib;
bifibrant when both. Replacement functors give $QX \xrightarrow{\sim} X
\xrightarrow{\sim} RX$. The homotopy category between bifibrant objects
is $\mathcal{C}(X, Y)/\!\sim$.

\textbf{Two foundational examples.}\quad Kan--Quillen on $\mathbf{sSet}$
(fibrations $=$ Kan fibrations; fibrant $=$ Kan complexes); Quillen on
$\mathbf{Top}$ (fibrations $=$ Serre; cofibrant $=$ CW). The adjunction
$\lvert{-}\rvert \dashv \mathrm{Sing}$ is a Quillen equivalence —
the homotopy hypothesis.

\textbf{Joyal model structure preview.}\quad Same cofibrations as
Kan--Quillen; fibrations only RLP against \emph{inner} horns. Fibrant
objects: quasi-categories. The home of $\infty$-category theory (Chapter~43).
\end{summarybox}


% ═══════════════════════════════════════════════════════════════════════════
\section*{Formal Restatement}
\addcontentsline{toc}{section}{Formal Restatement}
% ═══════════════════════════════════════════════════════════════════════════

\begin{formalrestatement}
\textbf{Quillen's axioms.}\quad $(\mathcal{C}, \mathcal{W}, \mathcal{C}\mathrm{of},
\mathcal{F}\mathrm{ib})$ satisfying:
\begin{itemize}
  \item (M1) $2$-out-of-$3$ for $\mathcal{W}$.
  \item (M2) Retract closure for $\mathcal{W}, \mathcal{C}\mathrm{of}, \mathcal{F}\mathrm{ib}$.
  \item (M3) $\mathrm{Mor}(\mathcal{C}) = \mathcal{C}\mathrm{of} \circ (\mathcal{W} \cap \mathcal{F}\mathrm{ib}) = (\mathcal{W} \cap \mathcal{C}\mathrm{of}) \circ \mathcal{F}\mathrm{ib}$.
  \item (M4) $\mathcal{C}\mathrm{of} = \mathrm{LLP}(\mathcal{W} \cap \mathcal{F}\mathrm{ib})$; $\mathcal{F}\mathrm{ib} = \mathrm{RLP}(\mathcal{W} \cap \mathcal{C}\mathrm{of})$.
  \item (M5) $\mathcal{C}$ has all small (co)limits.
\end{itemize}

\medskip
\textbf{Small object argument.}\quad For $\mathcal{I}$ a small set of maps with
small domains, every $f$ factors functorially as $\mathrm{cell}(\mathcal{I})
\circ \mathrm{RLP}(\mathcal{I})$.

\medskip
\textbf{Replacement.}\quad $Q : \mathcal{C} \to \mathcal{C}_{\mathrm{cof}}$
with $QX \xrightarrow{\sim} X$; $R : \mathcal{C} \to \mathcal{C}_{\mathrm{fib}}$
with $X \xrightarrow{\sim} RX$. Both natural in $X$.

\medskip
\textbf{Homotopy category.}\quad $\mathrm{Ho}(\mathcal{C})(X, Y) = [QX, RY]$
where $[-,-]$ denotes left ($=$ right) homotopy classes; this is the
small-hom-class model for $\mathcal{C}[\mathcal{W}^{-1}]$.

\medskip
\textbf{Quillen-$\mathbf{sSet}$.}\quad
\begin{itemize}
  \item $\mathcal{W}$: realization weak homotopy equivalences.
  \item $\mathcal{C}\mathrm{of}$: monomorphisms; generators
        $\partial\Delta^n \hookrightarrow \Delta^n$.
  \item $\mathcal{F}\mathrm{ib}$: Kan fibrations; generators of
        $\mathcal{W} \cap \mathcal{C}\mathrm{of}$:
        $\Lambda^n_k \hookrightarrow \Delta^n$ for $0 \leq k \leq n$.
\end{itemize}

\medskip
\textbf{Quillen-$\mathbf{Top}$.}\quad
\begin{itemize}
  \item $\mathcal{W}$: weak homotopy equivalences.
  \item $\mathcal{C}\mathrm{of}$: retracts of relative cell complexes;
        generators $S^{n-1} \hookrightarrow D^n$.
  \item $\mathcal{F}\mathrm{ib}$: Serre fibrations; generators of
        $\mathcal{W} \cap \mathcal{C}\mathrm{of}$:
        $D^n \hookrightarrow D^n \times [0, 1]$.
\end{itemize}

\medskip
\textbf{Quillen equivalence.}\quad $\lvert{-}\rvert : \mathbf{sSet}
\rightleftarrows \mathbf{Top} : \mathrm{Sing}$ is a Quillen equivalence;
$\mathrm{Ho}(\mathbf{sSet}) \simeq \mathrm{Ho}(\mathbf{Top})$.
\end{formalrestatement}
